% !TeX root = ../papers/nc-p3.tex
\section{正次元の階数一切断と核方向の剛性}
\label{sec:finite-marked-ruling-pencil}
\begin{theorem}\label{thm:finite-marked-ruling-pencil}
定理\ref{thm:finite-marked-projection-symmetry}の二つの四次元切断が一致するとする。
$\Sigma$ が正次元なら、その全体がSegreの一方の族の被約直線であり、
どの印でも $T$ はスキーム上の反対位置を同定する。従ってその側のテンソル対称性が従う。
\end{theorem}
\begin{proof}
$E$ の非零の片側ベクトルでの評価はどちらからも階数2以上である。
よって階数一点の整曲線の一方の射影が直線内にあれば、その全体は
一つの非零片側縮約の核に入り、線型包のベクトル次元は2以下である。
両射影とも平面を張るなら、その関数体の中の三次元座標空間二つの積は次元5以上である。
実際、滑らかな一点の位数で各空間に三つの異なる位数の基底を取り、
二つの三元整数集合の和集合（対ごとの和全体）に五つ以上の異なる値があることを使う。
これはSegre像の線型包が $\mathbb P^3$ に入ることに反する。
従って直線があり、左右を必要に応じて交換して、固定像を持つ
$K=\langle u v_1^t,u v_2^t\rangle$ と書く。

補元 $C,D$ を取り、上線を像直線 $ku$ を法とした二行三列への射影とする。
全ての $0\ne l\in u^\perp$ で $l^t C,l^t D$ は独立である。これは任意の
非零な $s\bar C+t\bar D$ の階数が2であることとも同値である。
従って $K$ の外には階数一点がない。$uv^t\in\mathbb P(K)$ においてSegreに接する
摂動は $ku$ を法として像の階数が高々1なので、$E$ 内の全接方向も $K$ に入る。
各閉点で局所次元と接次元が1なので交叉スキームはこの被約直線である。
さらに $\bar C,\bar D$ の共通右核は零である。さもなくば三列を二列に下げた
二行二列鉛筆が全非零元で可逆になり、その行列式の零点に反する。

二次元鉛筆が左右核を持ち、一般元が階数2なら、制限した二行二列の行列式が
$\Sigma$ の長さ2の部分を切る。その線型包は $\mathbb P(K)$ そのものなので、
実際には鉛筆の全元が階数1となり矛盾する。従って鉛筆は $K$ のみである。
共通右核は $kr_0=\ker v_1^t\cap\ker v_2^t$、左方向は任意の $[l]\in\mathbb P(u^\perp)$ である。
この記述が相対的にも成り立つことを示す。二平面と左右の二つの核直線をパラメータとする
Grassmann多様体と射影空間の積で縮約の消失スキームを取る。
$(K,l,r_0)$ で二平面の一次変形を、基底 $uv_i^t$ に $s_i C+t_i D$ を加えるものとする。
左縮約の一次式は
\[
 \alpha v_i^t+s_i l^t C+t_i l^t D=0,\qquad \alpha=\dot l^t u .
\]
$\alpha\ne0$ なら二行 $(s_i,t_i)$ は独立である。右縮約の一次式を $ku$ を法として読むと
$(s_i\bar C+t_i\bar D)r_0=0$ なので $r_0$ が共通右核となって矛盾する。
従って $\alpha=0$、$s_i=t_i=0$、$\dot r$ は射影的に零である。
このパラメータスキームは各閉点で接次元1、台は上の射影直線なので正則であり、
その被約射影直線と等しい。よって任意のArtin基底上でも右方向と鉛筆は固定である。

定理\ref{thm:finite-marked-secant-pencil}の局所消去を右方向で用いる。
印以外の誤差は、位置3で印と $r_0$ を結ぶ固定直線 $L$ のイデアルに消される。
印の接次元が1以下なら同定理の一様化変数の議論もそのまま適用できる。
評価には $(T x_0,x_3)$ における辺 $23$ の関係を用いる。これは $U_0U_3$ に属し、
辺 $30$ の自己直交性と支持補題が位置3の直線について適用できる。
接次元2で例外接方向がなければ前節までの三乗の証明による。
例外方向がある場合、、位置3の微分が単射なので $r_0$ へ行く射影接方向はただ一つである。
定理\ref{thm:finite-marked-tangent-secants}の二変数行行列の余核はここでも単純な一点である。
実際そこの双対数上の二独立元があれば上の相対剛性が右方向を一定にするが、
射影接直線上で右方向の微分は非零である。従って同定理の二次の余核による
$\mathfrak m^2 I_L f=0$ の証明は全て右方向で成り立つ。
支持長の矛盾により評価は零となる。接次元3は生じない。
実際、辺 $01$ の二つの端の評価像 $B,C$ は六次元空間の三平面で、
接次元は $\dim(B\cap C)$ である。辺 $12$ の消失から $w(p_0,-,p_2,-)$ の
行列は $B$ と $C^\perp$ の対合（順序は端の取り方による）であり、その階数は
$3-\dim(B\cap C)$。中支持4では $p_2$ によるこの行列は非零である。
\end{proof}
